% GENERATED FILE. Edit canonical lessons, then run npm run generate:books.
% canonical-source-hash: fnv1a64:160ec2d95bc7c6c8
% canonical-lessons: JA-C120-katai, JA-C120-tadashii, JA-C120-hima, JA-C120-shinsetsu, JA-C120-teinei, JA-R120-first-pass-animals-the-body-and-the-town, JA-R120-second-pass-shops-and-describing-words

\chapter{Hard, Correct, Free time, Kind, Polite}
\label{ch:ja-hard-correct-free-time-kind-polite}
\hlchaptermodality{120}

\begin{chapteropening}
\textbf{By the end of this chapter:} \emph{I can say hard, correct, free time, kind and polite.}

Complete the last lesson of chapter 120: I can say hard, correct, free time, kind and polite.
\end{chapteropening}

\section[\texorpdfstring{\ja{かたい} (katai)}{かたい (katai)}]{\ja{かたい} (katai) — hard}

\label{lesson:JA-C120-katai}



\begin{quote}
Before the new one: say the Japanese for a barber, then the Japanese for a butcher.
\end{quote}

\subsection*{You'll want to know: \ja{かたい}}
\textbf{\ja{かたい}} — \emph{katai} — \textquotedblleft{}hard\textquotedblright{}.

\begin{grammarlens}[title={What you've built}]
One more word for this chapter.
\end{grammarlens}

\subsection*{Your turn}
\begin{itemize}
\raggedright
  \item \emph{Say it:} \emph{katai}
  \item \emph{Say it:} \emph{katai}, once more
  \item \emph{Say it:} say \emph{katai}
  \item \emph{Recall:} say the Japanese for a barber, then the Japanese for a butcher, then say \emph{katai} again
\end{itemize}

\begin{tcolorbox}[breakable,colback=teal!4,colframe=teal!35!black,title={Before you move on}]
What does \ja{かたい} mean? (\textquotedblleft{}Hard\textquotedblright{}.) Say it once more.
\end{tcolorbox}
\section[\texorpdfstring{\ja{ただしい} (tadashii)}{ただしい (tadashii)}]{\ja{ただしい} (tadashii) — correct}

\label{lesson:JA-C120-tadashii}



\begin{quote}
Before the new one: say the Japanese for a butcher, then the Japanese for hard.
\end{quote}

\subsection*{You'll want to know: \ja{ただしい}}
\textbf{\ja{ただしい}} — \emph{tadashii} — \textquotedblleft{}correct\textquotedblright{}.

\begin{grammarlens}[title={What you've built}]
One more word for this chapter.
\end{grammarlens}

\subsection*{Your turn}
\begin{itemize}
\raggedright
  \item \emph{Say it:} \emph{tadashii}
  \item \emph{Say it:} \emph{tadashii}, once more
  \item \emph{Say it:} say \emph{tadashii}
  \item \emph{Recall:} say the Japanese for a butcher, then the Japanese for hard, then say \emph{tadashii} again
\end{itemize}

\begin{tcolorbox}[breakable,colback=teal!4,colframe=teal!35!black,title={Before you move on}]
What does \ja{ただしい} mean? (\textquotedblleft{}Correct\textquotedblright{}.) Say it once more.
\end{tcolorbox}
\section[\texorpdfstring{\ja{ひま} (hima)}{ひま (hima)}]{\ja{ひま} (hima) — free time; not busy}

\label{lesson:JA-C120-hima}



\begin{quote}
Before the new one: say the Japanese for hard, then the Japanese for correct.
\end{quote}

\subsection*{You'll want to know: \ja{ひま}}
\textbf{\ja{ひま}} — \emph{hima} — \textquotedblleft{}free time; not busy\textquotedblright{}.

\begin{grammarlens}[title={What you've built}]
One more word for this chapter.
\end{grammarlens}

\subsection*{Your turn}
\begin{itemize}
\raggedright
  \item \emph{Say it:} \emph{hima}
  \item \emph{Say it:} \emph{hima}, once more
  \item \emph{Say it:} say \emph{hima}
  \item \emph{Recall:} say the Japanese for hard, then the Japanese for correct, then say \emph{hima} again
\end{itemize}

\begin{tcolorbox}[breakable,colback=teal!4,colframe=teal!35!black,title={Before you move on}]
What does \ja{ひま} mean? (\textquotedblleft{}Free time; not busy\textquotedblright{}.) Say it once more.
\end{tcolorbox}
\section[\texorpdfstring{\ja{しんせつ} (shinsetsu)}{しんせつ (shinsetsu)}]{\ja{しんせつ} (shinsetsu) — kind}

\label{lesson:JA-C120-shinsetsu}



\begin{quote}
Before the new one: say the Japanese for correct, then the Japanese for free time; not busy.
\end{quote}

\subsection*{You'll want to know: \ja{しんせつ}}
\textbf{\ja{しんせつ}} — \emph{shinsetsu} — \textquotedblleft{}kind\textquotedblright{}.

\begin{grammarlens}[title={What you've built}]
One more word for this chapter.
\end{grammarlens}

\subsection*{Your turn}
\begin{itemize}
\raggedright
  \item \emph{Say it:} \emph{shinsetsu}
  \item \emph{Say it:} \emph{shinsetsu}, once more
  \item \emph{Say it:} say \emph{shinsetsu}
  \item \emph{Recall:} say the Japanese for correct, then the Japanese for free time; not busy, then say \emph{shinsetsu} again
\end{itemize}

\begin{tcolorbox}[breakable,colback=teal!4,colframe=teal!35!black,title={Before you move on}]
What does \ja{しんせつ} mean? (\textquotedblleft{}Kind\textquotedblright{}.) Say it once more.
\end{tcolorbox}
\section[\texorpdfstring{\ja{ていねい} (teinei)}{ていねい (teinei)}]{\ja{ていねい} (teinei) — polite, careful}

\label{lesson:JA-C120-teinei}



\begin{quote}
Before the new one: say the Japanese for free time; not busy, then the Japanese for kind.
\end{quote}

\subsection*{You'll want to know: \ja{ていねい}}
\textbf{\ja{ていねい}} — \emph{teinei} — \textquotedblleft{}polite, careful\textquotedblright{}.

\begin{grammarlens}[title={What you've built}]
One more word for this chapter.
\end{grammarlens}

\subsection*{Your turn}
\begin{itemize}
\raggedright
  \item \emph{Say it:} \emph{teinei}
  \item \emph{Say it:} \emph{teinei}, once more
  \item \emph{Say it:} say \emph{teinei}
  \item \emph{Recall:} say the Japanese for free time; not busy, then the Japanese for kind, then say \emph{teinei} again
\end{itemize}

\begin{tcolorbox}[breakable,colback=teal!4,colframe=teal!35!black,title={Before you move on}]
What does \ja{ていねい} mean? (\textquotedblleft{}Polite, careful\textquotedblright{}.) Say it once more.
\end{tcolorbox}
\section[(dialogue)]{Animals, the body and the town — a first pass}

\label{lesson:JA-R120-first-pass-animals-the-body-and-the-town}



\begin{quote}
Say the words for kind and polite.
\end{quote}

\subsection*{Your turn}
Say these aloud:

\begin{itemize}
\raggedright
  \item a farmhouse, a carpenter, a singer, work, a holiday
  \item a personal seal, a dumpling, soup stock, a bonito, wakame seaweed
  \item Chinese cabbage, garlic, sweet bean paste, honey, a net
  \item a lid, a target, a sunny spot, a sea cucumber, a sweetfish
\end{itemize}

\begin{tcolorbox}[breakable,colback=teal!4,colframe=teal!35!black,title={Before you move on}]
A farmhouse? (\textbf{\ja{のうか}}, \emph{nouka}.) A carpenter? (\textbf{\ja{だいく}}, \emph{daiku}.) A singer? (\textbf{\ja{かしゅ}}, \emph{kashu}.) Work? (\textbf{\ja{しごと}}, \emph{shigoto}.) A holiday? (\textbf{\ja{やすみ}}, \emph{yasumi}.) A personal seal? (\textbf{\ja{はんこ}}, \emph{hanko}.) A dumpling? (\textbf{\ja{だんご}}, \emph{dango}.) Soup stock? (\textbf{\ja{だし}}, \emph{dashi}.) A bonito? (\textbf{\ja{かつお}}, \emph{katsuo}.) Wakame seaweed? (\textbf{\ja{わかめ}}, \emph{wakame}.) Chinese cabbage? (\textbf{\ja{はくさい}}, \emph{hakusai}.) Garlic? (\textbf{\ja{にんにく}}, \emph{ninniku}.) Sweet bean paste? (\textbf{\ja{あんこ}}, \emph{anko}.) Honey? (\textbf{\ja{はちみつ}}, \emph{hachimitsu}.) A net? (\textbf{\ja{あみ}}, \emph{ami}.) A lid? (\textbf{\ja{ふた}}, \emph{futa}.) A target? (\textbf{\ja{まと}}, \emph{mato}.) A sunny spot? (\textbf{\ja{ひなた}}, \emph{hinata}.) A sea cucumber? (\textbf{\ja{なまこ}}, \emph{namako}.) A sweetfish? (\textbf{\ja{あゆ}}, \emph{ayu}.)
\end{tcolorbox}
\section[(dialogue)]{Shops and describing words — a second pass}

\label{lesson:JA-R120-second-pass-shops-and-describing-words}



\begin{quote}
Say the words for a trout and a crucian carp.
\end{quote}

\subsection*{Your turn}
Say these aloud:

\begin{itemize}
\raggedright
  \item a trout, a crucian carp, a young yellowtail, a Pacific saury, a short-neck clam
  \item a scallop, a flea, a bat, a Japanese serow, a water flask
  \item a rope, a ladder, a chain, cotton wool, a bookshop
  \item a florist, a fishmonger, a greengrocer, a barber, a butcher
  \item hard, correct, free time; not busy, kind, polite
\end{itemize}

\begin{tcolorbox}[breakable,colback=teal!4,colframe=teal!35!black,title={Before you move on}]
A trout? (\textbf{\ja{ます}}, \emph{masu}.) A crucian carp? (\textbf{\ja{ふな}}, \emph{funa}.) A young yellowtail? (\textbf{\ja{はまち}}, \emph{hamachi}.) A Pacific saury? (\textbf{\ja{さんま}}, \emph{sanma}.) A short-neck clam? (\textbf{\ja{あさり}}, \emph{asari}.) A scallop? (\textbf{\ja{ほたて}}, \emph{hotate}.) A flea? (\textbf{\ja{のみ}}, \emph{nomi}.) A bat? (\textbf{\ja{こうもり}}, \emph{koumori}.) A Japanese serow? (\textbf{\ja{かもしか}}, \emph{kamoshika}.) A water flask? (\textbf{\ja{すいとう}}, \emph{suitou}.) A rope? (\textbf{\ja{なわ}}, \emph{nawa}.) A ladder? (\textbf{\ja{はしご}}, \emph{hashigo}.) A chain? (\textbf{\ja{くさり}}, \emph{kusari}.) Cotton wool? (\textbf{\ja{わた}}, \emph{wata}.) A bookshop? (\textbf{\ja{ほんや}}, \emph{honya}.) A florist? (\textbf{\ja{はなや}}, \emph{hanaya}.) A fishmonger? (\textbf{\ja{さかなや}}, \emph{sakanaya}.) A greengrocer? (\textbf{\ja{やおや}}, \emph{yaoya}.) A barber? (\textbf{\ja{とこや}}, \emph{tokoya}.) A butcher? (\textbf{\ja{にくや}}, \emph{nikuya}.) Hard? (\textbf{\ja{かたい}}, \emph{katai}.) Correct? (\textbf{\ja{ただしい}}, \emph{tadashii}.) Free time; not busy? (\textbf{\ja{ひま}}, \emph{hima}.) Kind? (\textbf{\ja{しんせつ}}, \emph{shinsetsu}.) Polite? (\textbf{\ja{ていねい}}, \emph{teinei}.)
\end{tcolorbox}
